\begin{abstract}
Modern GPU-centric systems run heterogeneous workloads such as LLM inference, training, and vector search, whose performance and cost are dominated by policies for memory placement, kernel scheduling, and observability. Today, these policies are hard-coded inside framework-specific runtimes, proprietary vendor libraries, and monolithic kernel and device code, leaving operators unable to adapt them to workload mix, SLAs, or hardware changes without invasive code changes and disruption. Prior work has shown that externalizing CPU-side policies with eBPF---for scheduling, cache management, and congestion control---dramatically improves flexibility, but directly applying this pattern to GPUs is difficult. GPU driver mechanisms are tightly intertwined with low-level hardware interactions and lack safe extension points, and programmability on only the host or only the device cannot express coherent end-to-end policies.

We present \sys, a unified, safe, and efficient programmable policy runtime that spans both the CPU host and GPU devices. \sys refactors the GPU driver to expose stable, eBPF-based hooks at key control points, and offloads verified eBPF programs to execute directly inside GPU kernels via a lightweight, sandboxed PTX JIT, all over a shared map abstraction that provides cross-layer state. This design enables workload-specific policies for memory offloading, kernel admission and preemption, and fine-grained observability to run exactly where decisions occur, while preserving safety and low overhead.

We implement \sys on Linux with an extended \sysdriver and a GPU-side runtime, and evaluate it on realistic workloads including LLM inference, GNN training, vector search, and mixed-priority multi-tenant scenarios. Across these settings, \sys-based policies improve throughput by up to $5{\times}$, reduce tail latency by up to $2{\times}$, and add only a few percent runtime overhead for observability. These results demonstrate that unified CPU--GPU extensibility is a practical and necessary foundation for future GPU-centric systems.
\end{abstract}


\section{Introduction}

GPUs now dominate the compute and capital cost of modern datacenters, powering workloads that range from latency-sensitive LLM inference pipelines to large-scale training, graph analytics, and vector search. Across these workloads, end-to-end performance and efficiency are heavily determined not by raw FLOPs, but by \emph{policies}: how the system places and migrates data across device and host memory, how it admits and schedules kernels in multi-tenant environments, and how it collects observability signals to drive online adaptation. These policies must evolve as models, SLAs, and hardware change. Yet today they are largely frozen inside proprietary vendor stacks and per-framework runtimes, making GPU infrastructure rigid and costly to operate.

In current GPU stacks, the host driver and device firmware implement a fixed, closed set of mechanisms for memory management, context scheduling, and performance event collection. Frameworks such as PyTorch or TensorFlow, and emerging inference runtimes such as vLLM, build bespoke management logic on top of these mechanisms, often by defensively hoarding GPU memory or hard-wiring scheduling heuristics for a specific workload mix. Because the underlying mechanisms expose no safe extension points, operators cannot express workload-specific policies for memory placement or preemption at the points where they matter --- for example, when a kernel require a memory offloaded to CPU or when a low-priority job is about to delay a latency-critical query. As a result, datacenter GPU deployments suffer from unpredictable tail latency, underutilization, and frequent disruption when policies must be revised or tuned.

By contrast, the CPU side of modern systems has undergone a steady shift toward programmable policies. The extended Berkeley Packet Filter (eBPF) ecosystem now supports dynamically loaded, verified programs for scheduling (\texttt{sched\_ext}), page-cache eviction (\texttt{cache\_ext}), TCP congestion control (\texttt{struct\_ops}), and a wide body of observability tools~\cite{gregg16,gregg2021computing,Bachl21,Zhong22,jia2023programmable,bpftime}. In these designs, a small set of stable kernel mechanisms is exposed through structured attach points, and BPF programs implement policies that can be updated at runtime without rebooting or patching the kernel~\cite{schedext-docs,cacheext-docs}. It is tempting to apply the same idea to GPUs: treat the GPU driver as another kernel subsystem, attach host-side BPF programs to its control paths, and use those to implement adaptive GPU policies.

However, directly extending host-only BPF to GPU-centric systems fails for two fundamental reasons. First, many critical decisions --- such as whether to prefetch or evict certain data, how to throttle or prioritize individual warps, or when a kernel is triggering pathological memory behavior --- are made on the device, within GPU kernels and on-chip schedulers that are invisible to host-only instrumentation. Second, monolithic GPU drivers tightly couple resource management with low-level hardware protocols (e.g., MMU updates, interrupt handling, context teardown) and provide no stable, documented hooks; naively inserting probes risks GPU hangs, kernel panics, or subtle resource leaks. Conversely, GPU-only instrumentation tools such as CUPTI, NVBit, Neutrino, and eGPU~\cite{cupti,villa2019nvbit,neutrino,eGPU} provide rich telemetry by rewriting PTX or SASS, and recent userspace eBPF runtimes such as bpftime can inject BPF-like code into GPU kernels~\cite{bpftime}, but they are designed for profiling rather than online policy enforcement, lack a unified safety model with the OS, and do not integrate with privileged driver mechanisms.

Designing a general substrate for GPU policy programmability therefore raises several challenges. \emph{C1: Safe driver extensibility.} We must retrofit mechanism–policy separation into large, performance-critical GPU drivers without destabilizing their low-level interaction with hardware. \emph{C2: Safe and efficient device-side execution.} We need to execute user-supplied policy logic inside GPU kernels, on the critical path of memory accesses and synchronization, while guaranteeing bounded execution and memory safety in a massively parallel environment. \emph{C3: cross-layer state and metadata.} Policies often need to correlate host-side and device-side signals (e.g., page access patterns with warp-level stalls); doing so requires a shared state abstraction and efficient synchronization across the CPU–GPU boundary. \emph{C4: A single safety and control plane.} Finally, the policy runtime must support unified verification, deployment, and debugging for both host and device execution contexts; otherwise, operators are left with another fragmented stack of heterogeneous mechanisms.

We present \sys, a unified, safe, and efficient policy runtime for GPU-centric systems that addresses these challenges. \sys exposes a small set of stable, low-level GPU mechanisms---for memory placement and migration, kernel admission and preemption, and kernel-level instrumentation---as programmable attach points in an extended GPU driver. Policies are expressed as eBPF programs in a restricted instruction subset and are statically verified for memory safety and bounded execution. The same eBPF intermediate representation is then offloaded to GPU devices: \sys translates verified bytecode into PTX, injects it at kernel launch time via lightweight trampolines, and executes it inside kernels within a sandboxed runtime. Both host- and device-side policy code share a common map abstraction backed by unified memory, providing low-latency cross-layer state without manual data marshalling.

We implement \sys on Linux by extending \sysdriver over NVIDIA's open GPU kernel modules and by building a GPU-side PTX JIT and sandbox runtime. The driver modifications follow a \texttt{struct\_ops}-style pattern: they introduce a handful of attach points in memory-management and scheduling paths without restructuring existing mechanisms, and delegate decisions to verified eBPF callbacks. On the device side, lightweight PTX trampolines save minimal register state, invoke the translated eBPF snippet, and restore state to resume execution, adding sub-microsecond overhead in our microbenchmarks. Unlike prior GPU observability tools~\cite{cupti,villa2019nvbit,neutrino,eGPU} or userspace BPF runtimes~\cite{bpftime}, \sys offers one cross-layer policy model spanning the driver and the device, under a single verifier and control plane.

We evaluate \sys on realistic GPU datacenter workloads including LLM inference pipelines, graph neural network training, vector search, and mixed-priority multi-tenant scenarios. Using \sys, we implement adaptive memory placement and eviction, fine-grained kernel admission and preemption, and bcc-style observability tools for GPU kernels. Across these workloads, \sys-based policies improve throughput by up to $5{\times}$ and cut p99 latency by up to $2{\times}$ compared to static heuristics, while instrumentation-only deployments incur $2$–$3\%$ overhead. These results suggest that exposing unified, eBPF-style policy programmability across CPU and GPU is a practical and powerful abstraction for managing GPU-centric systems.

\begin{itemize}[leftmargin=*,itemsep=0pt]
    \item We identify and characterize key programmability gaps in GPU stacks for memory management, scheduling, and observability, and show how they lead to brittle, hard-coded policies that cannot adapt to evolving workloads.
    \item We introduce \sys, the first system to provide a unified eBPF-based policy runtime that spans CPU and GPU execution contexts, combining verified driver hooks with safe, JIT-compiled execution of BPF programs inside GPU kernels over a shared state abstraction.
    \item We demonstrate that \sys enables adaptive, workload-specific policies that significantly improve performance, utilization, and tail latency for realistic LLM inference, training, and vector workloads, with low overhead and without application or driver rewrites.
\end{itemize}
